% Existing published estimates only; no statistics were recomputed.
% Use \caption{\MainRawMainCaption} etc. within your paper's figure environment.
\newcommand{\MainRawMainCaption}{Main results on the frozen 500-source holdout at
$N=1024$ transmitted complex channel symbols, including the charged header.
Proposed (raw + partial + VAR), P1024, and SwinJSCC-80k (adapted) are compared in
PSNR, LPIPS-Alex, DINOv2-ViT-L/14 cosine similarity, and ConvNeXt prediction
agreement with the original source prediction (not classification accuracy).
Higher is better except for LPIPS. Means and shaded pointwise 95\% confidence
intervals are read directly from the published summary CSV. Each source first
averages its three frozen noise realizations; the existing 10,000-replicate
bootstrap weights the 500 sources equally, without multiplicity correction.
19 dB is outside Swin's training and calibration range. Its 19 dB marker is hollow
and its 13--19 dB segment is dashed. Swin uses six body channels (768 symbols) and
a 256-symbol protected header; this frozen adaptation and fixed 80k checkpoint
are not asserted to be the official optimum or fully converged. Equal source
count and total $N$ do not imply equal training, search, or header-protection
budgets. Only the three methods available in this holdout are plotted.}

\newcommand{\MainRawPartialCaption}{Partial-scale ablation on the same frozen
500-source holdout. Points, straight connecting lines, and error bars show the
published source-paired mean differences and pointwise 95\% confidence intervals
for PARTIAL $-$ WHOLE, with VAR completion used in both arms. The dashed horizontal
line marks zero. Positive differences favor PARTIAL for PSNR, DINOv2-L cosine
similarity, and ConvNeXt prediction agreement; negative differences favor PARTIAL
for LPIPS. ConvNeXt differences are percentage points, not relative percentages.
The intervals come directly from the paired CSV, not from subtraction of
single-method intervals. Both schemes have the same modulation and coding
permissions. PARTIAL can fall back to a whole-scale configuration: the final
WHOLE calibration winner remains eligible, and $K>0$ is not forced at every SNR.
The exact zero increments at 7 dB and the small or zero increments at 13 dB are
retained. Per-source three-noise means and the existing 500-source bootstrap are
used without recomputation or holdout reselection. This contrast is separate
from the VAR-completion ablation; the two sets of increments are not added.}

\newcommand{\MainRawCompletionCaption}{VAR-completion ablation on the frozen
500-source holdout, with the PARTIAL transmission configuration held fixed.
Points, straight connecting lines, and error bars show the published source-paired
mean differences and pointwise 95\% confidence intervals for VAR completion $-$
Direct decoding; the dashed horizontal line marks zero. Both arms use the same
actually received tokens from each transmission and the same frozen decoder
$D_c$. Direct decoding does not generate untransmitted residuals: their residual
embedding contribution is set to zero. This is not filling unknown tokens with
codebook index zero, and Direct receives no additional source-image ground truth.
Positive differences favor VAR completion for PSNR, DINOv2-L cosine similarity,
and ConvNeXt prediction agreement; negative differences favor VAR completion
for LPIPS. ConvNeXt differences and their intervals are expressed in percentage
points. All estimates and intervals are read from the existing paired CSV after
the original three-noise source averaging and 500-source bootstrap; no new
bootstrap is performed. This contrast is not added to PARTIAL $-$ WHOLE.}
